\honest{Every Cause Wants To Be A Cult}{Eliezer Yudkowsky}{2007}

The Register alleges that Wikipedia's top administrators use a secret mailing list to ban critics, that one of them banned a productive editor as a spy for a critics' site, solely because the editor was so productive, and that the top people closed ranks. \nb{The same article reports that the ban was lifted ``after a mere 75 minutes,'' that the administrator gave up her powers, and that the Arbitration Committee admonished her. The post does not mention this.}

Is there something wrong with the cause of organizing the world's knowledge? ``Correspondence bias alert!'' Ordinary human nature explains it. In-group bias, happy death spirals and spirals of hate are all ordinary. A noble cause needs no hidden flaw to turn cultish; ``It is sufficient that the adherents be human.'' One true idea does not switch off the halo effect, or status games, or in-group bias.

As heat evens out and programs turn into patches, ``every Cause wants to be a cult.'' Every group with an unusual goal ``will trend toward the cult attractor unless they make a constant effort to resist it,'' as a house stays cool only while the air conditioner runs. \nb{The evidence for ``every'' is the Wikipedia case and the Objectivists. The post does not say what counts as the effort, so it names no case that could count against the claim. Five days earlier, in ``Evaporative Cooling of Group Beliefs,'' a cooler group was a more fanatical one.}

A group whose cry was ``Rationality! Reason! Objective reality!'' went semicultish. Calling your cause rationality protects you no more than a sign saying ``Cold!'' Worshipping rationality will not make you sane any more than worshipping gravity lets you fly. You can use probability theory, but you cannot join it.

Cultishness is ``quantitative, not qualitative.'' Even science has its battle lines. Do journals favor famous authors, or unknown authors from famous institutions? How well does blind review work? I choose this example over the vague complaint that scientists are not open to new ideas because it shows a place where the effort is actually being made. \nb{Peters and Ceci had tested the first question in 1982. Published articles resubmitted under fictitious names and institutions were mostly rejected: ``eight of the nine'' that were not recognized. The post asks the question and does not cite the answer.}

This essay is not a catalog of techniques. My point is only that a worthy cause needs as much effort as any other, and that naming its battle lines does not betray it. Asking ``Cultish, yes or no?'' is like sorting engines into perfectly efficient and inefficient instead of measuring the waste. If you think cults are made of ``mutants,'' you will not make the effort.
